\documentclass{article}
\usepackage[utf8]{inputenc}
\usepackage[T1]{fontenc}
\usepackage[swedish]{babel}
\usepackage{graphicx}
\usepackage{enumitem}
\usepackage{fancyhdr}
\usepackage[table]{xcolor}
\usepackage{tabularx}
\usepackage{changepage}
\usepackage[most]{tcolorbox}
\usepackage{amsmath}
\usepackage{csquotes}
\usepackage{hyperref}
\usepackage{url}

\hypersetup{
    colorlinks=true,
    linkcolor=blue,
    urlcolor=blue,
    pdftitle={Metodbeskrivning för riskbedömning av kund v2.0}
}
\title{Metodbeskrivning för riskbedömning av kund\\\large Version 2.0 - Integrerad med Celestial Risk Engine}
\pagestyle{fancy}
\fancyhf{}
\fancyhead[R]{Metodbeskrivning för riskbedömning av kund v2.0, \today}

\begin{document}

\vspace{1em}

\maketitle
\tableofcontents
\newpage

\section*{Inledning}
\addcontentsline{toc}{section}{Inledning}

Denna metodbeskrivning fastställer hur vi som redovisningsbyrå genomför riskbedömningen av kunder, både vid antagandet av nya uppdrag och vid den löpande uppdateringen av befintliga kunders riskprofiler. Syftet är att säkerställa en enhetlig och rättssäker tillämpning av penningtvättslagstiftningens krav på kundkännedom och riskklassificering.

\textbf{Nytt i version 2.0:} Denna uppdaterade version integrerar vår klassiska PTL-metodik (Hot × Sårbarhet) med automatiserad datainsamling via Roaring API samt analys av bokföringsunderlag (SIE-filer) genom vår Celestial Risk Engine. Detta skapar en tvådelad bedömning:

\begin{enumerate}
\item \textbf{PTL-metodik (juridisk grund):} Hot × Sårbarhet = KYC-nivå (standard eller skärpt)
\item \textbf{Celestial Engine (operationell):} 0--100 poäng för intern prioritering och beslut
\end{enumerate}

För att upprätthålla ett aktuellt och rättssäkert arbetssätt uppdateras detta dokument minst en gång per år, eller vid behov till följd av ändrad lagstiftning, vägledning från tillsynsmyndigheter eller interna riskobservationer. Ansvarig för uppdatering och godkännande är den verksamhetsansvarige. Vid varje ändring uppdateras dokumentets versionsnummer.

% Definitioner: hot och sårbarhet
\section*{Definition av risknivåer}
\addcontentsline{toc}{section}{Definition av risknivåer}

För att kunna uppfylla penningtvättslagens krav på riskbaserade åtgärder (2 kap. 1 § och 3 §) krävs en systematisk och dokumenterad metod för att identifiera och bedöma risker. 

Vi inleder därför med en definition av de risknivåer som tillämpas i vår verksamhet. Risknivån utgör en kombination av externa hot (kundrisk) och interna sårbarheter (tjänstens känslighet), vilket ligger till grund för beslut om kontrollåtgärder och kundacceptans.

\begin{itemize}[label=--]
  \item \textbf{Hot} – den yttre risken: i vilken grad en viss kundtyp eller verksamhet kan utnyttjas för penningtvätt eller finansiering av terrorism. Bedömningen utgår från den nationella riskbedömningen och avser hur attraktiv kunden är för kriminella aktörer.

  \item \textbf{Sårbarhet} – den inre risken: i vilken utsträckning våra egna tjänster kan manipuleras eller missbrukas utan att det upptäcks. Faktorer inkluderar dokumentation, kontrollrutiner, spårbarhet och insyn.
\end{itemize}

\subsection*{Princip för riskbedömning}
\addcontentsline{toc}{subsection}{Princip för riskbedömning}

Vi tillämpar en enkel och transparent modell där:

\begin{quote}
\textbf{Risknivå = Hot (kund) × Sårbarhet (tjänst)}
\end{quote}

Hotnivån hämtas från den \textit{Nationella riskbedömningen av penningtvätt och finansiering av terrorism i Sverige 2020/2021}\footnote{\url{https://www.ekobrottsmyndigheten.se/wp-content/uploads/2021/05/nationell-riskbedomning-av-penningtvatt-och-finansiering-av-terrorism-2020-2021.pdf}}. Sårbarhetsnivån fastställs internt för varje tjänst utifrån dess kontrollbarhet.

\textbf{Syftet är att styra vilka åtgärder som krävs enligt penningtvättslagen, med riskpoängen som beslutsunderlag.}

%%%%%%%%%%%%%%%%%%%%%%%%%%%

\vspace{3em}

%% Allmänt Hot x Sårbarhet

\section*{Metod för allmän riskbedömning}
\addcontentsline{toc}{section}{Metod för allmän riskbedömning}

Vår metod för allmän riskbedömning baseras på den riskmatris och skala som presenteras i \textit{Nationell riskbedömning av penningtvätt och finansiering av terrorism 2020/2021}, utgiven av bland andra Finansdepartementet, Polismyndigheten, Ekobrottsmyndigheten och Finansinspektionen.

Risknivån beräknas genom att multiplicera en bedömd hotnivå med en bedömd sårbarhetsnivå enligt följande formel:

\[
\text{Risknivå}_{\text{PTL}} = \text{Hotnivå} \times \text{Sårbarhetsnivå}
\]

Varje komponent bedöms på en fyrgradig skala:

\begin{itemize}[label=--]
  \item 1 = Låg
  \item 2 = Medel
  \item 3 = Betydande
  \item 4 = Hög
\end{itemize}

\noindent Skalan följer den semikvantitativa metodik som används i den nationella riskbedömningen (bilaga "Process och metod", s. 123 ff.). Genom att använda denna metod säkerställs att vår klassificering är spårbar, jämförbar och förankrad i myndighetsgemensamma riktlinjer.

För varje kombination av kundtyp och tjänst motiverar vi bedömda siffror i separata tabeller i efterföljande avsnitt.

%%%%%%%%%%%%%%%%%%%%%%%%%%%%%
\vspace{3em}

\subsection*{Härledning av hotnivåer per verksamhetsområde enligt myndighetskatalog}
\addcontentsline{toc}{subsection}{Härledning av hotnivåer per verksamhetsområde enligt myndighetskatalog}

Som underlag för vår generella riskbedömning har vi utgått från den myndighetsgemensamma rapporten \textit{Nationell riskbedömning av penningtvätt och finansiering av terrorism i Sverige 2020/2021}, utgiven av Ekobrottsmyndigheten, Polismyndigheten, Finansinspektionen m.fl.

Dokumentet innehåller i kapitel 7 (s. 39–117) en sektorskatalog där samtliga verksamhetsområden som omfattas av penningtvättslagstiftningen bedöms individuellt utifrån risken att utnyttjas för penningtvätt eller finansiering av terrorism. För varje sektor anges en risknivå i termerna \textit{låg}, \textit{medel}, \textit{betydande} eller \textit{hög}, med tillhörande motivering.

\subsubsection*{Automatisk mappning från SNI-kod till hotnivå}

\textbf{Nytt i v2.0:} Vi hämtar företagens branschklassificering automatiskt via Roaring API, som returnerar SNI-kod (5 siffror). Denna kod mappas sedan till myndighetens hotnivå enligt följande tabell:

\begin{center}
\small
\begin{tabular}{|p{7cm}|c|c|}
\hline
\textbf{Verksamhetsområde} & \textbf{SNI-kod} & \textbf{Hotnivå} \\
\hline
\multicolumn{3}{|c|}{\textbf{HÖG RISK (4)}} \\
\hline
Bygg och anläggning & 41.xxx, 42.xxx, 43.xxx & 4 \\
Restaurang och catering & 56.xxx & 4 \\
Städ och fastighetsservice & 81.xxx & 4 \\
Landtransport & 49.xxx & 4 \\
Varuhandel (kontantintensiv) & 47.xxx & 4 \\
Bolagsbildare och företagsmäklare & 70.xxx & 4 \\
Spelverksamhet & 92.xxx & 4 \\
\hline
\multicolumn{3}{|c|}{\textbf{BETYDANDE RISK (3)}} \\
\hline
Fastighetsmäklare & 68.3xx & 3 \\
Bokförings-, revisions- och skatterådgivning & 69.2xx & 3 \\
Oberoende jurister (ej advokatbyråer) & 69.10x (ej 69.10.1) & 3 \\
Kontorshotell och postboxföretag & 82.11x & 3 \\
Finansiell verksamhet (bank, holdingbolag) & 64.1xx, 64.2xx, 64.3xx & 3 \\
Elektroniska pengar och betalningsinstitut & 66.xxx & 3 \\
Försäkring och pension & 65.xxx & 3 \\
\hline
\multicolumn{3}{|c|}{\textbf{MEDEL RISK (2)}} \\
\hline
IT-konsulter och programutveckling & 62.xxx, 63.xxx & 2 \\
Arkitekter, ingenjörer, teknisk konsult & 71.xxx & 2 \\
Reklam och marknadsföring & 73.xxx & 2 \\
Advokater och advokatbyråer & 69.10.1 & 2 \\
Valutaväxling och övrig finansiell verksamhet & 64.9xx & 2 \\
Forskning och utveckling & 72.xxx & 2 \\
\hline
\multicolumn{3}{|c|}{\textbf{LÅG RISK (1)}} \\
\hline
Offentlig förvaltning & 84.xxx & 1 \\
Utbildning (offentlig) & 85.xxx & 1 \\
Hälso- och sjukvård (offentlig) & 86.xxx & 1 \\
Kommunala bolag & Ägarstruktur & 1 \\
\hline
\end{tabular}
\end{center}

\textbf{Tillämpning:} När en ny kund registreras hämtas SNI-kod från Roaring API. Systemet matchar automatiskt mot ovanstående tabell och tilldelar korrekt hotnivå (1--4). Om SNI-koden inte finns i tabellen tilldelas standardvärdet \textbf{2 (Medel risk)}.

\subsubsection*{Översättning av hotnivåer till numerisk skala}

För att möjliggöra en enhetlig och kvantitativ tillämpning av riskbedömningen har vi översatt myndigheternas verbala riskkategorier till en numerisk skala enligt följande:

\begin{center}
\begin{tabular}{|c|c|}
\hline
\textbf{Benämning i myndighetsrapport} & \textbf{Numeriskt värde i vår modell} \\
\hline
Låg & 1 \\
Medel & 2 \\
Betydande & 3 \\
Hög & 4 \\
\hline
\end{tabular}
\end{center}

Vi tillämpar denna fyrgradiga skala både för hotnivå (extern exponering) och sårbarhetsnivå (intern kontrollförmåga). Den resulterande risknivån erhålls genom multiplikation:

\[
\text{Riskvärde}_{\text{PTL}} = \text{Hotnivå} \times \text{Sårbarhetsnivå}
\]

Detta ger ett slutligt intervall från 1 till 16 och ligger till grund för våra prioriteringar och åtgärdsplaner enligt penningtvättslagen.

\begin{tcolorbox}
\textbf{Exempel:} En restaurangverksamhet (SNI 56.101) har hotnivå 4 enligt myndighetens klassificering. Våra redovisningstjänster har sårbarhetsnivå 3. Detta ger ett PTL-riskvärde på \(4 \times 3 = 12\), vilket klassificeras som \textit{hög risk} och kräver skärpt kundkännedom.
\end{tcolorbox}

%%%
\newpage
\section*{Celestial Risk Engine - Operationell riskbedömning}
\addcontentsline{toc}{section}{Celestial Risk Engine - Operationell riskbedömning}

\subsection*{Syfte och avgränsning}

\textbf{Nytt i v2.0:} Förutom den juridiskt förankrade PTL-metoden tillämpar vi en kompletterande \textit{Celestial Risk Engine} för att få en mer detaljerad och operationell riskbild. Denna metod kombinerar:

\begin{enumerate}
\item \textbf{15 Roaring Risk Indicators} - Offentliga data från Bolagsverket, konkurser, PEP, styrelsebyten etc.
\item \textbf{22 Bokföringsflaggor} - Analys av SIE-filer för att upptäcka fraud-mönster
\item \textbf{10 Cross-validation rules} - Integritetskontroller mellan Roaring-data och bokföring
\end{enumerate}

\textbf{Viktigt:} Celestial Engine \textit{ersätter inte} PTL-metoden. De två metoderna har olika syften:

\begin{center}
\begin{tabular}{|p{6cm}|p{6cm}|}
\hline
\textbf{PTL-metodik} & \textbf{Celestial Engine} \\
\hline
Juridisk grund för KYC-nivå & Operationell prioritering \\
Output: Låg/Normal/Hög risk & Output: 0--100 poäng \\
Beslut: Standard vs Skärpt KYC & Beslut: Accept/Enhanced DD/Reject \\
Baseras på bransch (Hot × 3) & Baseras på faktisk verksamhet \\
\hline
\end{tabular}
\end{center}

\subsection*{Beräkningsmodell}

Celestial Engine ger varje kund en poäng mellan 0 och 100 enligt följande viktade formel:

\[
\text{Score}_{\text{Celestial}} = (S_{\text{Roaring}} \times W_{\text{public}}) + (S_{\text{Bokföring}} \times W_{\text{accounting}}) + (S_{\text{Validation}} \times W_{\text{cross}})
\]

\noindent där:
\begin{itemize}
\item $S_{\text{Roaring}}$ = Poäng från Roaring Risk Indicators (0--100)
\item $S_{\text{Bokföring}}$ = Poäng från 22 bokföringsflaggor (0--100)
\item $S_{\text{Validation}}$ = Poäng från cross-validation (0--100)
\item $W_{\text{public}}$, $W_{\text{accounting}}$, $W_{\text{cross}}$ = Vikter som summerar till 1.0
\end{itemize}

\subsection*{Dynamisk viktning baserat på PTL Hot-nivå}

\textbf{Nyckelfunktion:} Viktningen anpassas automatiskt baserat på kundens Hot-nivå från PTL-metoden:

\begin{center}
\begin{tabular}{|c|c|c|c|l|}
\hline
\textbf{Hot} & \textbf{$W_{\text{public}}$} & \textbf{$W_{\text{accounting}}$} & \textbf{$W_{\text{cross}}$} & \textbf{Motivering} \\
\hline
4 (Hög) & 0.20 & 0.60 & 0.20 & \parbox{5cm}{Bygg/restaurang: offentlig data mindre tillförlitlig, fokus på bokföring} \\
\hline
3 (Betydande) & 0.25 & 0.55 & 0.20 & Standard balance \\
\hline
2 (Medel) & 0.25 & 0.55 & 0.20 & Standard balance \\
\hline
1 (Låg) & 0.35 & 0.45 & 0.20 & \parbox{5cm}{Offentlig sektor: offentlig data mycket tillförlitlig} \\
\hline
\end{tabular}
\end{center}

\textbf{Rationalen:} I högriskverk samheter (bygg, restaurang) arbetar företag ofta svart, vilket gör offentlig data (Roaring) mindre tillförlitlig. Därför ökar vi vikten på bokföringsanalys till 60\%. I offentlig sektor är det tvärtom - offentlig data är mycket tillförlitlig, så vi ökar vikten för Roaring till 35\%.

\subsection*{Roaring Risk Indicators (15 regler)}

Via Roaring API hämtas följande 15 indikatorer automatiskt:

\begin{enumerate}
\item \textbf{fTaxReg} - F-skatteregistrering (false = alarm)
\item \textbf{vatReg} - Momsregistrering (false = alarm om omsättning >80k)
\item \textbf{arbAvgReg} - Arbetsgivaravgift (false = alarm om anställda >0)
\item \textbf{registrationDate} - Nyetablerad (<6 månader = alarm)
\item \textbf{boardChangeCount} - Styrelsebyten (>3/år = instabilitet)
\item \textbf{addressChangeDates} - Adressbyten (>3/år = misstänkt)
\item \textbf{industryCodeChanges} - Branschbyten (>2/år = verksamhetsändring)
\item \textbf{orgDataChanges} - Organisationsdata ändringar (>3/3 år)
\item \textbf{bankruptciesPerReference} - Personliga konkurser (>1 = red flag)
\item \textbf{connectedBankruptcyCompanies} - Nätverk konkurser (>5 = varning)
\item \textbf{pepCount} - PEP-status (>3 = skärpt KYC)
\item \textbf{legalCount} - Rättsliga åtgärder (>3 = dispyter)
\item \textbf{revenueByEmployee} - Omsättning/anställd (<1000 kr = fraud)
\item \textbf{missingAuditor} - Saknar revisor (true = varning om >3M omsättning)
\item \textbf{meanAgeOfRepresentatives} - Medelålder styrelse (35--50 = optimalt)
\end{enumerate}

Varje indikator som triggar alarm ger 5--20 poäng beroende på allvarlighetsgrad. Totalen omvandlas till en skala 0--100.

\subsection*{Bokföringsflaggor (22 regler)}

Genom analys av kundens SIE-fil identifieras följande 22 varningsflaggor:

\begin{enumerate}
\item Fakturor med jämna belopp (>30\% = round-tripping)
\item Bruten nummerordning (missing invoices = oredovisad försäljning)
\item Kontantinsättningar >22,000 SEK (>2000 EUR = obligatorisk skärpt KYC)
\item Oförklarliga utbetalningar (>50k utan faktura)
\item Betalning före förfallodatum utan rabatt (>5 st = samordning)
\item Stora aktieägartillskott (>500k = kontrollera ursprung)
\item Penningflöde till högriskländer
\item Frekventa ägarbyten (>2/år)
\item Privatbostadskostnader i bokföringen
\item Oförklarliga in- och utbetalningar
\item Motpart i högriskland
\item Stora internationella transaktioner utan tydligt syfte
\item Fakturor med varierande utseende (olika format = flera företag?)
\item Snabba ägarbyten vid platsbesök
\item Verksamhet vid platsbesök överensstämmer ej med bokföring
\item Kontinuerligt stora förluster (>3 år = oklar finansiering)
\item Negativt eget kapital (konkursrisk)
\item Likvida medel <1 månads burn rate
\item Överkrediterad (kontokredit >limit)
\item Leverantörsskulder >90 dagar
\item Omsättning/anställd avviker från branschgenomsnitt (>200\%)
\item Bruttomarginal <branschmedian -50\%
\end{enumerate}

Varje flagga ger 5--50 poäng beroende på severity (Låg/Medel/Hög/Kritisk). Totalen omvandlas till 0--100.

\subsection*{Cross-validation (10 regler)}

Dessa regler jämför Roaring-data mot bokföringen för att hitta diskrepanser:

\begin{enumerate}
\item Omsättning i SIE vs Roaring (>20\% skillnad = false reporting)
\item Antal anställda i SIE (via lönekostnad) vs Roaring (>2 skillnad = fake employees)
\item VAT-registrering i Roaring vs moms i SIE (false = fraud)
\item Arbetsgivaravgift i Roaring vs SIE konto 2730 (mismatch = fraud)
\item F-skatt i Roaring vs faktureringsaktivitet i SIE
\item Styrelsebyten i Roaring vs signaturändringar i SIE
\item PEP-status vs högrisk transaktioner i SIE
\item Registreringsdatum i Roaring vs första SIE-post (>6 mån = backdating)
\item Revisor saknas i Roaring men >10M omsättning i SIE (illegal enligt ÅRL)
\item Adressändringar i Roaring vs adressdata i SIE
\end{enumerate}

Varje discrepans ger 10--60 poäng. Totalen omvandlas till 0--100.

%%%
\newpage
\subsection*{Exempel: Restaurang med fraud-mönster}

\textbf{Företag:} Mama Mia Pizzeria AB (SNI 56.101)

\textbf{Steg 1: PTL-metodik (juridisk grund)}

\begin{itemize}
\item Hot-nivå: 4 (Restaurang = Hög risk enligt myndigheten)
\item Sårbarhet: 3 (Redovisningstjänst = Betydande)
\item PTL-riskvärde: $4 \times 3 = 12$
\item Klassificering: \textbf{Hög risk}
\item KYC-krav: \textbf{Skärpt kundkännedom (3 kap 18 § PTL)}
\end{itemize}

\textbf{Steg 2: Celestial Engine (operationell)}

\textit{Roaring Indicators (2 alarms):}
\begin{itemize}
\item arbAvgReg = false (10 poäng)
\item missingAuditor = true (5 poäng)
\item \textbf{Roaring Score:} 15/100 (låg public risk)
\end{itemize}

\textit{Bokföringsflaggor (2 kritiska):}
\begin{itemize}
\item 5 kontantinsättningar >22,000 SEK (50 poäng - KRITISK)
\item Bruten fakturaserie \#45--\#67, 23 st saknas (30 poäng - HÖG)
\item \textbf{Bokföring Score:} 80/100 (hög fraud risk)
\end{itemize}

\textit{Cross-validation (OK):}
\begin{itemize}
\item Omsättning: SIE 3.2M vs Roaring 3.0M = 6.7\% skillnad (OK)
\item \textbf{Validation Score:} 0/100 (good consistency)
\end{itemize}

\textit{Viktning (Hot 4):}
\[
\text{Score}_{\text{Celestial}} = 15 \times 0.20 + 80 \times 0.60 + 0 \times 0.20 = 3 + 48 + 0 = 51
\]

Justerat för Hot 4 bransch: $51 + 10 = 61/100$

\textbf{Steg 3: Beslut}

\begin{tcolorbox}[colback=red!10]
\textbf{PTL-beslut:} Skärpt KYC (Hot 4 bransch)

\textbf{Celestial-beslut:} ACCEPT med Enhanced DD (score 61)

\textbf{Kombinerat:} ACCEPTERA kund MEN:
\begin{itemize}
\item OBLIGATORISK skärpt kundkännedom (5 kap 3 § 2 st PTL - kontant >22k)
\item Begär förklaring på alla 5 kontantinsättningarna
\item Begär förklaring på saknade fakturor \#45--\#67
\item Kvartalsvis uppföljning (högriskverksamhet)
\item Platsbesök inom 6 månader
\item Alla kontanttransaktioner >22k dokumenteras separat
\end{itemize}
\end{tcolorbox}

%%%
\newpage
\section*{Motivering av sårbarhetsnivåer per tjänst}
\addcontentsline{toc}{section}{Motivering av sårbarhetsnivåer per tjänst}

Vi väljer i denna analys att klassificera samtliga våra tjänster med sårbarhetsnivån \textbf{3 – betydande}, i linje med slutsatserna i den myndighetsgemensamma rapporten från 2020/2021. Tjänsterna bygger i stor utsträckning på kundens egna uppgifter och är därmed exponerade för manipulation, även om våra interna rutiner syftar till att motverka detta. 

Vi tillämpar denna nivå som en försiktighetsåtgärd och är öppna för revidering vid förändrade förutsättningar eller vägledning från tillsynsmyndighet.

\subsection*{1. Löpande bokföring – Sårbarhet = 3 (betydande)}
Enligt den nationella riskbedömningen (kap. 7.15, s. 91) är redovisningsbyråer en särskilt utsatt sektor eftersom de i praktiken kan möjliggöra eller maskera brott genom manipulativ bokföring. Tjänsten bygger på kundens underlag och är ofta beroende av kundens ärlighet. Vi bedömer därför sårbarheten som \textbf{betydande}.

\textbf{Nytt i v2.0:} Genom automatisk analys av SIE-filer (22 bokföringsflaggor) minskar vi dock denna sårbarhet i praktiken, eftersom manipulation upptäcks systematiskt.

\subsection*{2. Momsdeklaration – Sårbarhet = 3 (betydande)}
Momsfusk utgör enligt rapporten ett vanligt och välkänt modus för ekonomisk brottslighet (s. 40--41). Även om deklarationen i sig är formbunden, är underlaget mottagligt för manipulation. Därför klassificeras sårbarheten som \textbf{betydande}.

\subsection*{3. Årsbokslut/Årsredovisning – Sårbarhet = 3 (betydande)}
Årsbokslut och årsredovisningar är till stor del sammanställningar av redan hanterade transaktioner, men enligt rapporten (s. 85) används de i hög grad för att vilseleda banker och myndigheter. Detta innebär en betydande sårbarhet eftersom manipulation ofta sker genom selektiv presentation av data.

\subsection*{4. Inkomstdeklaration – Sårbarhet = 3 (betydande)}
Trots att deklarationen skickas till Skatteverket sker själva datainsamlingen internt eller via konsult. Deklarationen är därför beroende av kundens rapportering och dokumentation, vilket öppnar för otillbörlig påverkan. Riskbedömningen stöder att sårbarheten för denna typ av tjänst bedöms som \textbf{betydande}.

%%%
\newpage
\section*{Sammanfattande beslutsmodell}
\addcontentsline{toc}{section}{Sammanfattande beslutsmodell}

Vår tvådelade riskbedömning resulterar i följande beslutsmatris:

\begin{center}
\renewcommand{\arraystretch}{1.8}
\begin{tabular}{|c|c|c|p{6cm}|}
\hline
\textbf{PTL-risk} & \textbf{Celestial} & \textbf{Beslut} & \textbf{Åtgärder} \\
\hline
Låg (1--3) & 0--30 & Accept & Standard KYC, årlig uppföljning \\
\hline
Låg (1--3) & 31--60 & Accept & Standard KYC, halvårsvis uppföljning \\
\hline
Låg (1--3) & 61--80 & Enhanced DD & Fördjupad granskning innan accept \\
\hline
Låg (1--3) & 81--100 & Reject & Fraud bevis trots låg PTL-risk \\
\hline
Normal (4--6) & 0--30 & Accept & Standard KYC, halvårsvis uppföljning \\
\hline
Normal (4--6) & 31--60 & Accept & Standard KYC, kvartalsvis uppföljning \\
\hline
Normal (4--6) & 61--80 & Enhanced DD & Fördjupad kontroll + extra dokumentation \\
\hline
Normal (4--6) & 81--100 & Reject & Fraud upptäckt \\
\hline
Hög (7+) & 0--30 & Accept & Skärpt KYC, kvartalsvis uppföljning \\
\hline
Hög (7+) & 31--60 & Accept & Skärpt KYC, månatlig uppföljning \\
\hline
Hög (7+) & 61--80 & Enhanced DD & Skärpt KYC + VD-godkännande \\
\hline
Hög (7+) & 81--100 & Reject & Kombinerad hög risk + fraud \\
\hline
\end{tabular}
\end{center}

\subsection*{Särskilda krav vid kontanthantering}

Enligt 5 kap 3 § 2 st penningtvättslagen är skärpt kundkännedom \textbf{obligatorisk} vid:

\begin{itemize}
\item Enstaka transaktion motsvarande minst 2000 euro (ca 22,000 SEK) i kontanter
\item Flera sammanhängande transaktioner som tillsammans uppgår till minst 2000 euro
\end{itemize}

\textbf{Celestial Engine flaggar automatiskt} alla kontantinsättningar >22,000 SEK i SIE-filen och kräver skärpt KYC oavsett övrig riskprofil.

%%%
\newpage
\section*{Dokumentation och spårbarhet}
\addcontentsline{toc}{section}{Dokumentation och spårbarhet}

Varje riskbedömning ska dokumenteras och arkiveras enligt följande:

\subsection*{1. PTL-riskbedömning (juridisk dokumentation)}
\begin{itemize}
\item Hot-nivå (1--4) med källhänvisning till myndighetsrapport
\item Sårbarhetsnivå (konstant 3 för redovisningstjänster)
\item Riskvärde (Hot × 3)
\item Klassificering: Låg/Normal/Hög risk
\item KYC-nivå: Standard eller Skärpt
\item Eventuella justeringar (PEP +4, högriskland +3, nyetablerad +2 etc.)
\item Beslutsdatum och ansvarig
\end{itemize}

\subsection*{2. Celestial Engine rapport (operationell dokumentation)}
\begin{itemize}
\item Roaring Score (0--100) med lista på triggade indikatorer
\item Bokföring Score (0--100) med lista på flaggor
\item Validation Score (0--100) med lista på discrepanser
\item Total Celestial Score efter viktning
\item Top 5 riskfaktorer
\item Beslut: Accept/Enhanced DD/Reject
\item Alla Roaring API-svar arkiverade (JSON)
\item SIE-fil analyserad (arkiverad separat)
\end{itemize}

\subsection*{3. Kombinerat beslutsdokument}
Ett sammanfattande dokument som inkluderar:
\begin{itemize}
\item Båda riskbedömningarna (PTL + Celestial)
\item Slutligt beslut med motivering
\item Specificerade åtgärder och uppföljningsplan
\item Signatur från verksamhetsansvarig
\end{itemize}

Dokumentationen ska bevaras i minst 5 år från det att affärsförbindelsen upphört (4 kap 8 § penningtvättslagen).

%%%
\newpage
\section*{Uppdatering och revidering}
\addcontentsline{toc}{section}{Uppdatering och revidering}

Denna metodbeskrivning ska ses över och vid behov uppdateras:

\begin{itemize}
\item Minst en gång per år
\item Vid ändrad lagstiftning eller nya föreskrifter
\item Vid ny vägledning från Länsstyrelsen eller annan tillsynsmyndighet
\item Vid väsentliga förändringar i den nationella riskbedömningen
\item Vid identifierade brister i den egna verksamheten
\item Vid uppdatering av Roaring API eller Celestial Engine
\end{itemize}

\textbf{Version 2.0 (2025-10-24):} Integration av Celestial Risk Engine med automatiserad datainsamling via Roaring API och bokföringsanalys. Separation mellan PTL-metodik (juridisk) och operationell riskbedömning (Celestial). Dynamisk viktning baserat på Hot-nivå. Uppdatering till SNI 2025.

\textbf{Tidigare versioner:}
\begin{itemize}
\item Version 1.0: Grundläggande PTL-metodik (Hot × Sårbarhet)
\end{itemize}

\vspace{2em}

\noindent
\rule{\textwidth}{0.4pt}

\vspace{1em}

\noindent
\textbf{Godkänd av:} \_\_\_\_\_\_\_\_\_\_\_\_\_\_\_\_\_\_\_\_\_\_\_\_\_\_

\noindent
\textbf{Datum:} \_\_\_\_\_\_\_\_\_\_\_\_\_\_\_\_\_\_\_\_\_\_\_\_\_\_

\noindent
\textbf{Nästa revidering:} \_\_\_\_\_\_\_\_\_\_\_\_\_\_\_\_\_\_\_\_\_\_\_\_\_\_

\end{document}
